% Phan cua thanh vien B: (d) Heuristic va tim kiem cuc bo, va muc Tabu Search trong (e).
% Day la tep con (khong co preamble). Can cac goi: amsmath, algorithm, algpseudocode, booktabs.
% Cac o [CHO SO LIEU] dien sau khi chay thuc nghiem; khong dien so chua do.

\section{Heuristic xây dựng và tìm kiếm cục bộ}

\subsection{Biểu diễn lời giải và cấu trúc dữ liệu}
Lời giải là tập cột $X\subseteq\{1,\dots,n\}$, lưu bằng mảng boolean $sel[j]$. Với mỗi hàng $i$ ta duy trì
$cnt[i]=|\{j\in X: i\in S_j\}|$ (số cột đang chọn phủ hàng $i$). Lời giải hợp lệ khi và chỉ khi $cnt[i]\ge 1$ với mọi $i$.
Thêm hoặc bỏ cột $j$ cập nhật $cnt$ trong $O(|S_j|)$ và cập nhật chi phí trong $O(1)$.
Bộ đọc dựng cả danh sách theo hàng ($R_i$: các cột phủ hàng $i$) và theo cột ($S_j$).

\subsection{Tham lam Chvátal}
Tại mỗi bước chọn cột chưa chọn có tỉ số $c_j/|S_j\cap U|$ nhỏ nhất, với $U$ là tập hàng chưa được phủ.

\begin{algorithm}[htbp]
\caption{Tham lam Chvátal}
\begin{algorithmic}[1]
\State $X\gets\emptyset$; $U\gets\{1,\dots,m\}$
\While{$U\neq\emptyset$}
  \State $j^*\gets\arg\min_{j\notin X,\,|S_j\cap U|>0}\; c_j/|S_j\cap U|$
  \State $X\gets X\cup\{j^*\}$; $U\gets U\setminus S_{j^*}$
\EndWhile
\State \Call{XoaCotThua}{$X$}
\State \Return $X$
\end{algorithmic}
\end{algorithm}

Chỉ xét các cột chứa ít nhất một hàng thuộc $U$. Mỗi vòng lặp tốn $O(\sum_{j}|S_j|)$ trong trường hợp xấu nhất, và có tối đa $m$ vòng.
Chvátal (1979) chứng minh tỉ số xấp xỉ $H(d)\le 1+\ln d$ với $d=\max_j|S_j|$.

\paragraph{Xóa cột thừa.} Duyệt các cột đã chọn theo chi phí giảm dần; cột $j$ bị bỏ nếu $cnt[i]\ge 2$ với mọi $i\in S_j$
(bỏ đi vẫn hợp lệ). Độ phức tạp $O(\sum_{j\in X}|S_j|)$.
Trên ví dụ nhỏ của đề, tham lam cho $\{4,5,2,1\}$ chi phí 19; bước xóa cột thừa loại cột 5 và đưa về 16 (tối ưu).

\subsection{Thủ tục sửa chữa}
\textsc{SuaChua} thêm cột theo đúng tiêu chí tỉ số của Chvátal cho đến khi phủ hết các hàng còn thiếu; có tham số $forbid$ (tập cột bị cấm thêm).
Nếu mọi ứng viên đều bị cấm, bỏ lệnh cấm để bảo đảm lời giải hợp lệ. Các ứng viên cùng tỉ số được chọn ngẫu nhiên theo seed.

\subsection{Hai cấu trúc lân cận}
\textbf{$N_1$: hoán đổi 1--1.} Bỏ cột $j\in X$ và thêm cột $k\notin X$. Gọi $T_j=\{i\in S_j: cnt[i]=1\}$ là các hàng chỉ do $j$ phủ.
Phép hoán đổi hợp lệ khi và chỉ khi $T_j\subseteq S_k$. Mức thay đổi chi phí:
\[
\Delta = c_k - c_j .
\]
Nhờ mảng $cnt$, không phải tính lại toàn bộ lời giải: tính $T_j$ tốn $O(|S_j|)$, các ứng viên $k$ chỉ lấy từ danh sách cột của một hàng thuộc $T_j$.
Nếu $T_j=\emptyset$ thì $j$ là cột thừa và bị bỏ.

\textbf{$N_2$: bỏ cột rồi sửa chữa.} Bỏ $j$, gọi \textsc{SuaChua} với $forbid=\{j\}$, rồi xóa cột thừa.
Mức thay đổi $\Delta=$ (chi phí mới) $-$ (chi phí cũ), tính từ giá trị chi phí được cập nhật gia tăng. Lân cận này cho phép thay một cột bằng nhiều cột rẻ hơn.

\subsection{Leo đồi: first improvement và best improvement}
Duyệt các cột đã chọn theo thứ tự ngẫu nhiên (theo seed).
\begin{itemize}
  \item \textbf{First improvement (\texttt{HC\_FI}):} áp dụng ngay phép di chuyển đầu tiên có $\Delta<0$.
  \item \textbf{Best improvement (\texttt{HC\_BI}):} xét mọi $j\in X$ ở cả hai lân cận, áp dụng phép có $\Delta$ nhỏ nhất.
\end{itemize}
Dừng khi không còn phép di chuyển cải thiện hoặc hết ngân sách thời gian.
Lời giải xuất phát là kết quả của tham lam Chvátal, nên mỗi lần chạy không tệ hơn lời giải xuất phát của chính nó.

\paragraph{Chi phí một bước.} $N_1$: $O(|S_j|)$ để tìm $T_j$ cộng chi phí kiểm tra từng ứng viên $k$.
$N_2$: dominated bởi \textsc{SuaChua}, tức $O(\sum_j |S_j|)$ mỗi lần; vì vậy $N_2$ đắt hơn $N_1$.
Không dùng delta (tính lại toàn bộ) thì mỗi ứng viên tốn thêm $O(\sum_{i}|R_i|)$ để kiểm tra phủ.

\paragraph{So sánh FI và BI.} Kết quả thực nghiệm trên 10 thể hiện (4 nhỏ tự sinh, 3 vừa tự sinh, 3 chuẩn OR-Library \texttt{scp41}, \texttt{scp51}, \texttt{scp61}) với 10 seed (\texttt{1..10}), ngân sách $0{,}5$ giây được tổng hợp trong bảng sau (chạy bằng \texttt{experiments/member\_b/run\_member\_b.py}, lưu tại \texttt{results/member\_b/eval\_member\_b.csv}):
\begin{center}
\small
\begin{tabular}{lrrrrrr}
\toprule
Thể hiện (Ref) & GREEDY & HC\_FI Mean (Gap\%) & $T_{\mathrm{FI}}$ (s) & HC\_BI Mean (Gap\%) & $T_{\mathrm{BI}}$ (s) \\
\midrule
\texttt{nho\_10x20\_s101} (65)  & 65  & 65,00 (0,00\%)  & 0,0004 & 65,00 (0,00\%)  & 0,0003 \\
\texttt{nho\_15x25\_s102} (136) & 136 & 136,00 (0,00\%) & 0,0003 & 136,00 (0,00\%) & 0,0007 \\
\texttt{nho\_20x28\_s103} (212) & 216 & 212,00 (0,00\%) & 0,0004 & 212,00 (0,00\%) & 0,0008 \\
\texttt{nho\_25x30\_s104} (133) & 133 & 133,00 (0,00\%) & 0,0004 & 133,00 (0,00\%) & 0,0004 \\
\texttt{vua\_100x400\_s201} (211) & 213 & 211,00 (0,00\%) & 0,0121 & 211,00 (0,00\%) & 0,0142 \\
\texttt{vua\_150x600\_s202} (261) & 272 & 269,10 (3,10\%) & 0,0274 & 268,20 (2,76\%) & 0,0302 \\
\texttt{vua\_200x1000\_s203} (316) & 330 & 316,50 (0,16\%) & 0,0800 & 316,00 (0,00\%) & 0,1277 \\
\texttt{scp41} (429) & 434 & 433,00 (0,93\%) & 0,0549 & 433,00 (0,93\%) & 0,0594 \\
\texttt{scp51} (253) & 269 & 268,00 (5,93\%) & 0,1011 & 268,00 (5,93\%) & 0,1120 \\
\texttt{scp61} (138) & 142 & 141,00 (2,17\%) & 0,0560 & 141,00 (2,17\%) & 0,0684 \\
\bottomrule
\end{tabular}
\end{center}
Cả hai biến thể đều cải thiện rõ rệt so với \texttt{GREEDY} (ví dụ trên \texttt{vua\_200x1000\_s203} giảm từ 330 xuống 316). \texttt{HC\_BI} ổn định hơn nhẹ ở các thể hiện vừa (\texttt{vua\_150x600\_s202}, \texttt{vua\_200x1000\_s203}) nhưng tốn nhiều thời gian quét mỗi vòng hơn \texttt{HC\_FI} (ví dụ $0{,}1277$s so với $0{,}0800$s trên \texttt{vua\_200x1000\_s203}). Trên \texttt{scp51} và \texttt{scp61}, cả hai bị kẹt tại cực tiểu cục bộ ($268$ và $141$), cho thấy cần metaheuristic để thoát khỏi cực tiểu cục bộ.

\section{Tabu Search (phần trong mục các metaheuristic)}

\subsection{Nguồn gốc}
Tìm kiếm Tabu do Glover (1986) đề xuất: luôn đi đến lân cận tốt nhất không bị cấm, dùng danh sách tabu để chặn quay lại và tiêu chí mong đợi (aspiration).

\subsection{Thiết kế cho bài toán phủ tập}
\begin{itemize}
  \item \textbf{Biểu diễn:} tập cột hợp lệ (như trên); \emph{luôn giữ hợp lệ} bằng sửa chữa, không dùng hàm phạt.
  \item \textbf{Lời giải đầu:} tham lam Chvátal có ngẫu nhiên hóa khi hòa.
  \item \textbf{Lân cận:} bỏ một cột $j\in X$ rồi sửa chữa và xóa cột thừa (lân cận $N_2$). Mỗi vòng xét tối đa 40 cột chọn ngẫu nhiên từ $X$.
  \item \textbf{Thuộc tính tabu:} cột $j$ vừa bị bỏ không được thêm lại trong $tenure+r$ vòng, với $r$ ngẫu nhiên trong $\{0,1,2\}$.
  \item \textbf{Độ dài tabu mặc định:} $tenure=\max(5,\lfloor 0.1|X|\rfloor+3)$.
  \item \textbf{Aspiration:} một phép di chuyển có dùng cột đang bị cấm vẫn được chấp nhận nếu cho chi phí nhỏ hơn lời giải tốt nhất đã biết $c(X^*)$; ngược lại phép di chuyển được tính lại với các cột tabu bị cấm. Nếu không còn ứng viên nào khác để phủ hàng thiếu thì nới cấm để giữ tính hợp lệ.
  \item \textbf{Tiêu chí dừng:} hết ngân sách thời gian. Mọi lời giải báo cáo đều qua bộ kiểm tra độc lập của A.
\end{itemize}

\begin{algorithm}[htbp]
\caption{Tabu Search cho Set Cover}
\begin{algorithmic}[1]
\State $X\gets$ \Call{Chvatal}{ }; $X^*\gets X$; $it\gets 0$
\While{còn thời gian}
  \State $it\gets it+1$; $best\gets\emptyset$
  \For{$j\in$ mẫu tối đa 40 cột của $X$}
    \State $X'\gets$ \Call{SuaChua}{$X\setminus\{j\}$, $forbid=\{j\}$}; xóa cột thừa
    \If{$X'$ dùng cột tabu \textbf{và} $c(X')\ge c(X^*)$}
      \State $X'\gets$ \Call{SuaChua}{$X\setminus\{j\}$, $forbid=\{k: tabu[k]>it\}\cup\{j\}$}; xóa cột thừa
    \EndIf
    \State giữ $X'$ nếu tốt hơn $best$
  \EndFor
  \State $X\gets best$ (kể cả khi tệ hơn $X$); $tabu[j]\gets it+tenure+r$
  \If{$c(X)<c(X^*)$} $X^*\gets X$ \EndIf
\EndWhile
\State \Return $X^*$
\end{algorithmic}
\end{algorithm}

\subsection{Độ phức tạp mỗi vòng lặp}
Mỗi vòng thử tối đa 40 phép bỏ--sửa chữa (một số phép phải sửa chữa lần hai khi dùng cột tabu), mỗi phép $O(\sum_j|S_j|)$ trong trường hợp xấu nhất.
Phép thử được hoàn tác bằng nhật ký thêm/bỏ cột (rollback), chỉ tốn công tỉ lệ với số cột bị thay đổi, không sao chép cả lời giải.

\subsection{Ý nghĩa tham số và phân tích độ nhạy}
$tenure$ điều khiển cân bằng khai thác và khám phá: nhỏ thì dễ quay lại vùng cũ (kẹt), lớn thì cấm quá nhiều cột tốt.
Kích thước mẫu (40) đánh đổi giữa chất lượng lân cận và số vòng lặp trong cùng ngân sách.
Để tránh dùng chính tập đánh giá cuối, việc tinh chỉnh được thực hiện trên tập riêng gồm hai thể hiện tự sinh bởi \texttt{src/gen\_scp.py}: \texttt{tune\_100x400\_s901} (seed 901) và \texttt{tune\_150x600\_s902} (seed 902), mỗi cấu hình chạy 10 seed (\texttt{1..10}) với ngân sách $0{,}3$ giây (kết quả chi tiết trong \texttt{results/member\_b/ts\_sensitivity.csv}):

\begin{center}
\small
\begin{tabular}{lcrrrr}
\toprule
Cấu hình (\texttt{tenure}; \texttt{sample}) & \multicolumn{2}{c}{\texttt{tune\_100x400\_s901}} & \multicolumn{2}{c}{\texttt{tune\_150x600\_s902}} & Evals TB (\texttt{s902}) \\
 & Best & Mean $\pm$ Std & Best & Mean $\pm$ Std & \\
\midrule
\texttt{tenure=3; sample=40}    & 274 & 274,00 $\pm$ 0,00 & 296 & 296,00 $\pm$ 0,00 & 1339,7 \\
\texttt{tenure=5; sample=40}    & 274 & 274,00 $\pm$ 0,00 & 296 & 296,00 $\pm$ 0,00 & 1427,8 \\
\texttt{tenure=10; sample=40}   & 274 & 274,00 $\pm$ 0,00 & 285 & 289,50 $\pm$ 4,77 & 1473,9 \\
\texttt{tenure=15; sample=40}   & 274 & 274,00 $\pm$ 0,00 & \textbf{284} & \textbf{284,00 $\pm$ 0,00} & 1398,3 \\
\texttt{tenure=auto; sample=40} & 274 & 274,00 $\pm$ 0,00 & 285 & 293,30 $\pm$ 4,40 & 1491,2 \\
\texttt{tenure=auto; sample=10} & 274 & 274,10 $\pm$ 0,32 & \textbf{284} & 288,00 $\pm$ 4,37 & 1532,2 \\
\texttt{tenure=auto; sample=20} & 274 & 274,00 $\pm$ 0,00 & 285 & 292,40 $\pm$ 5,56 & 1312,2 \\
\texttt{tenure=auto; sample=80} & 274 & 274,00 $\pm$ 0,00 & 285 & 293,50 $\pm$ 4,50 & 1721,4 \\
\bottomrule
\end{tabular}
\end{center}
Trên \texttt{tune\_150x600\_s902}, \texttt{tenure} nhỏ ($3, 5$) bị kẹt tại $296$ do vòng lặp quay lại trạng thái cũ; khi tăng \texttt{tenure} lên $10$--$15$, thuật toán thoát cực tiểu cục bộ và đạt $284{,}00$ ($\mathrm{std}=0{,}00$ tại \texttt{tenure=15}). Với ngân sách ngắn $0{,}3$ giây, \texttt{sample} nhỏ ($10$) giúp thực hiện nhiều vòng lặp hơn mỗi giây nên đạt trung bình tốt hơn ($288{,}00$ so với $293{,}50$ của \texttt{sample=80}), trong khi ở các thể hiện chuẩn cỡ lớn (\texttt{scp51}), \texttt{TS} mặc định vẫn hạ chi phí từ $268$ (\texttt{HC}) xuống tốt nhất $256$ (trung bình $261{,}40$).
